(1) 连续性 \(d\rho/\rho+du/u+dA/A=0\)；定常 Euler 方程 \(u\,du+dp/\rho=0\)，等熵声速关系 \(dp=c^2d\rho\)，于是 \(d\rho/\rho=-M^2du/u\)。消去密度变化得到 \(dA/A=(M^2-1)du/u\)。它显示亚声速加速需收缩，超声速加速需扩张，并要求喉部过音速时 \(dA=0\)。

(2) 绝热焓守恒 \(c_pT+u^2/2=c_pT_0\)，其中 \(c_p=\gamma R/(\gamma-1)\)，从而 \(T/T_0=[1+(\gamma-1)M^2/2]^{-1}\)。等熵 \(p/\rho^\gamma=\mathrm{const}\) 联合理想气体式得到 \(p/p_0=[1+(\gamma-1)M^2/2]^{-\gamma/(\gamma-1)}\)。静温静压与总状态不可混淆。

(3) 用 \(\rho u=pM\sqrt{\gamma/(RT)}\) 消去静态变量，得到 \(\dot m/A=p_0\sqrt{\gamma/(RT_0)}M[1+(\gamma-1)M^2/2]^{-(\gamma+1)/(2(\gamma-1))}\)。对 \(M\) 求导为零的唯一正根是 \(M=1\)，二阶变化为负。最大质量通量即 \(p_0\sqrt{\gamma/(RT_0)}[2/(\gamma+1)]^{(\gamma+1)/(2(\gamma-1))}\)。

(4) 在壅塞态喉部声速，收缩区使亚声速气流加速，扩张区才能使已超声速气流继续加速。降低背压并不无限增加固定喉面积的通量。若喷管内有激波，静态压强和温度会突跃、熵增加、总压下降；这违反题设全程可逆等熵，必须另解激波前后边界匹配。
